\subsection{Held-out heterogeneous and translation fields}
The confirmatory heterogeneous campaign recovered 23,570 of 24,641 visible windows (95.65\%), with no false accepted visible windows. The paired mean gain over the archived baseline was 23.65 percentage points, with a 95\% pair-bootstrap interval of 14.78--35.19 points. Against a matched re-solve control, the mean gain was 23.40 points (14.69--34.61). Figure~\ref{fig:paired} shows why the pooled count is insufficient: difficult pairs account for substantial recovery, and a small aggregate difference can conceal a local failure.

The translation campaign recovered 28,641 of 29,459 visible windows, compared with 25,173 for the corrected baseline, with no pair regression and no visible false acceptance. An earlier comparison had rejected near-identity affine prefits and was not used as the confirmatory baseline. Retaining that correction is essential because pure translation is a basic control for the geometric preprocessing itself.

\begin{figure}[tb]\centering\includegraphics[width=.92\linewidth]{figures/paired_coverage.pdf}\caption{Paired heterogeneous-field coverage reconstructed from the archived confirmatory report. Lines connect the baseline and propagated result on each image pair, sorted by gain. Both methods use that pair's visible requested centres as denominator. These are final correlation outcomes, not direct feature-match counts.}\label{fig:paired}\end{figure}

\subsection{Independent same-centre transfer}
The five-family experiment changes support size while preserving centres. Table~\ref{tab:multi} and Figure~\ref{fig:multi} show increased coverage in every pooled family, with no loss of previously correct centres and no visible false acceptance in these draws. Nevertheless, gradient error increased on the common correct population in every family. In heterogeneous fields, displacement root-mean-square error decreased from 0.0589 to 0.0191 pixels, while gradient root-mean-square error increased from 0.0009 to 0.0025. Smaller support can reduce non-affine mismatch while containing less information for derivative estimation.

\begin{table}[tb]\centering\caption{Independent transfer at identical centres, five pairs per family. Correct means an accepted solve with displacement error below two pixels.}\label{tab:multi}
\begin{tabular}{lrrr}\toprule Family & Visible & 64-pixel correct & 32-pixel correct\\\midrule
Translation&6467&6322&6395\\Rigid rotation&6278&6008&6144\\Affine&5845&5616&5718\\Heterogeneous&6061&5804&5920\\Discontinuous&6556&6414&6443\\\bottomrule\end{tabular}\end{table}
\begin{figure}[tb]\centering\includegraphics[width=\linewidth]{figures/multifamily.pdf}\caption{Coverage after smaller-support transfer at the same requested centres. Bars start at zero to preserve the scale of improvement. No visible false acceptance occurs in these draws; this does not establish universal safety. The severe-shear qualification remains failed.}\label{fig:multi}\end{figure}

Figures~\ref{fig:spatial} and~\ref{fig:discontinuous} connect the window-level statistics to the actual speckle images. Each shows the full image pair, the spatial outcome, and native-pixel supports before and after the archived local affine correction. The local example is the valid correct centre nearest the image centre, excluding interface-crossing supports; it is not selected by the smallest photometric residual. The pullback uses the estimated six parameters, never the analytical displacement. Full-frame views are reduced for printing, whereas the enlarged supports retain native sampling. The absolute intensity difference is a visual diagnostic, not the normalized acceptance criterion or a displacement-error map. No gain or offset is fitted for these displays; illumination changes can therefore leave a non-zero intensity difference despite successful geometric registration.

\begin{figure}[p]\centering\includegraphics[width=\linewidth]{figures/spatial_transfer.pdf}
\caption{Synthetic heterogeneous pair 62051005/62091005, at $3840\times2160$ pixels. (a,b) Actual reference and deformed speckle images. (c) Archived outcome at the requested centres: blue is correct at both support sizes, green is newly recovered at 32 pixels, and orange is unresolved or incorrect among visible centres. (d) A native 32-pixel reference support, located by the magenta cross. (e) The target support centred at the estimated corresponding position, with translation correction only. (f) The target pulled back using the complete estimated local affine transformation. (g) Absolute intensity difference between (d) and (f), on a fixed 0--0.25 scale; larger differences saturate. The same grayscale range is used for (d)--(f). No image is warped with analytical truth. These are final solver outputs after initialization and transfer, not initializer-only precision.}\label{fig:spatial}\end{figure}

\begin{figure}[p]\centering\includegraphics[width=\linewidth]{figures/discontinuous_speckle.pdf}
\caption{Synthetic discontinuous pair 62051004/62091004. Panels, scales and outcome definitions match Figure~\ref{fig:spatial}. The full images show the separated material regions; no smooth global remapping is used to bridge the interface. The enlarged 32-pixel support belongs to one continuous region and therefore does not demonstrate successful correlation across a crack. This archived family has a grid-aligned opening, a limitation of the experiment.}\label{fig:discontinuous}\end{figure}
\subsection{Recorded stages and propagation provenance}
Figure~\ref{fig:chainexample} follows one existing heterogeneous pair through the recorded proposal and solve states. The pre-correlation affine state and the post-correlation initial state are archived separately. The plot does not reconstruct an unavailable raw score volume or claim that a projected centre was directly matched. The illustration is a secondary archive analysis, not an additional validation pair.

\begin{figure}[p]\centering\includegraphics[width=\linewidth]{figures/recorded_chain.pdf}
\caption{Recorded stages for heterogeneous pair 62051005/62091005. (a) Thirty-two evenly indexed hard candidates from the archived 1296 pairs, shown as arrows on the native reference; arrows are proposals, not verified matches. (b) Native 64-pixel reference support at the solver-accepted centre nearest the image centre. (c) Target pullback with archived affine-only initial parameters. (d) Target pullback with archived affine-plus-local-correlation initial parameters evaluated at that centre; this is the interpolated proposal, not a reconstructed raw correlation maximum. (e) Pullback after the archived native IC-GN solve. (f) Absolute final intensity difference. All transformations use saved estimated states, never truth. This is the 64-pixel initial solve, distinct from the 32-pixel support-transfer illustration.}\label{fig:chainexample}\end{figure}

\begin{figure}[p]\centering\includegraphics[width=\linewidth]{figures/propagation_limits.pdf}
\caption{Secondary archive analysis of propagation and its limits. Top: heterogeneous pair 24051003/24091003, with nine confirmed anchors. Bottom: shear pair 59051005/59091005 at peak gradient 2.0. Left: confirmed-anchor locations; middle: recorded accepted generations; right: final analytical scoring (strict error below two native pixels), with rejected visible centres and false acceptances distinguished by symbols. The severe-shear row contains 22 false acceptances; numerical acceptance did not establish correctness. Exact-centre truth is used only for this scoring panel. These selected cases illustrate mechanisms and are not estimates of family-wide performance.}\label{fig:propagation}\end{figure}

\subsection{Failure envelope and experimental images}
At peak shear gradient 2.0, the frozen severe-shear analysis recorded 57 incorrect visible acceptances among 4,080 accepted windows (1.40\%). All 57 were also incorrect with analytical local-affine initialization. Thus the local model, support or solver limits performance as well as the initializer. Better seed positions alone cannot be assumed to remove this failure. Numerical acceptance is not a guarantee of physical correctness.

A separately frozen image-residual suspicion rule detected 119 of 120 false accepted coarse states on an additional heterogeneous and extreme-shear sample, but also flagged 221 of 6,979 correct states. Guarded smaller-support retries recovered 7,070 correct and one false result among 8,075 visible shear centres, leaving 199 previously correct flagged centres unresolved. This is a tradeoff, not a universal error detector. An unknown discontinuity is especially difficult because a nominal neighbourhood can combine two material motions.

Figure~\ref{fig:experimental} shows the experimental pair supplied by Adrien (frames 000000 and 000075), recorded at $7920\times6004$ pixels. Of 7,280 requested centres valid under the archived region-of-interest mask, 7,097 were accepted by the solver (97.49\%). This denominator excludes masked-out centres from the 11,439-centre rectangular grid. The enlargement is the accepted centre nearest the image centre, selected without ranking photometric residuals. The 64-pixel pullback uses its archived final affine parameters. This real-image example uses the masked propagation campaign, whereas Figures~\ref{fig:spatial}--\ref{fig:discontinuous} use the separate 32-pixel transfer campaign; they are not a matched accuracy comparison.

\begin{figure}[p]\centering\includegraphics[width=\linewidth]{figures/experimental_speckle.pdf}
\caption{Experimental images and local remapping, Adrien frames 000000--000075. (a,b) Native-image fields of view, reduced for print. (c) Green marks solver-accepted centres inside the supplied mask; orange marks masked-domain centres not accepted. Unmarked regions are not assessed by this map. There is no independent displacement truth. (d--f) A 64-pixel reference support, target support with translation correction only, and target support with the complete estimated local affine correction. (g) Absolute intensity difference, with the same fixed scale as the synthetic examples. Agreement of speckle patterns demonstrates visual registration at this location, not calibrated displacement or strain accuracy. No global field interpolation or oracle transformation is used for this enlargement.}\label{fig:experimental}\end{figure}

The Adrien experimental pair provides masked-image integration evidence. Without independent truth, accepted counts and visual remapping cannot establish displacement accuracy. We therefore do not combine those counts with synthetic correct-coverage statistics. Definitive experimental qualification requires acquisition calibration, documented speckle statistics, and independent displacement or strain evidence.

\subsection{Performance evidence}
Component time and complete scripted time are different quantities. On one synthetic native-resolution pair, an isolated paired comparison reduced median script wall time from 16.451 to 11.541 seconds by caching solver pixels. Three measured repetitions followed one warm-up and alternated arm order. Both arms obtained 1,060 correct of 1,121 visible windows, with identical states, status and provenance. These times include image generation, process startup, diagnostics and writing results; they are not kernel latency or camera throughput. The GPU was not exclusively reserved, although no competing benchmark was observed at the recorded checks.

A four-pair independent check of shared prefit and deferred rotation generation preserved 5,011 correct of 5,242 visible windows, identical acceptance masks and zero visible false acceptances. No additional rotation hypotheses were requested on these pairs. This establishes parity for that sample, not the cost of rotational fallback. Current measurements do not support a subsecond end-to-end claim. Future timing must include transfers, failed hypotheses and required result recovery, with the same images and window budgets.

